%% ========================================================================
%% Qin Jiushao — Shuxue Jiuzhang (Mathematical Treatise in Nine Sections)
%% Fascicles 7–9: Bilingual Classical Chinese → Modern Chinese Edition
%% ========================================================================
%% Engine: XeLaTeX
%% Layout: Parallel columns (LEFT: Classical Chinese, RIGHT: Modern Chinese)
%% ========================================================================
\documentclass[11pt,a4paper]{article}

%% -------- Core Packages --------
\usepackage{fontspec}
\usepackage{polyglossia}
\usepackage{amsmath,amssymb}
\usepackage{geometry}
\usepackage{graphicx}
\usepackage{xcolor}
\usepackage{booktabs}
\usepackage{longtable}
\usepackage{xeCJK}
\usepackage{paracol}
\usepackage{ulem}

%% -------- Page Geometry --------
\geometry{
  paperwidth=210mm,paperheight=297mm,
  left=18mm,right=18mm,top=25mm,bottom=25mm,
  headheight=14pt,footskip=30pt
}

%% -------- Column Ratio --------
\columnratio{0.50,0.50}

%% -------- Language Setup --------
\setmainlanguage{english}

%% -------- Font Configuration --------
\setmainfont{LinLibertine_R.otf}[
  BoldFont=LinLibertine_RB.otf,
  ItalicFont=LinLibertine_RI.otf,
  BoldItalicFont=LinLibertine_RBI.otf
]
\setsansfont{LinBiolinum_R.otf}[
  BoldFont=LinBiolinum_RB.otf,
  ItalicFont=LinBiolinum_RI.otf
]
\setmonofont{DejaVu Sans Mono}[Scale=MatchLowercase]

%% Chinese Fonts (xeCJK) — using Noto Sans CJK SC
\setCJKmainfont[Path=/mnt/agents/fonts/noto-sc-extracted/]{NotoSansCJKsc-Regular.otf}
\setCJKsansfont[Path=/mnt/agents/fonts/noto-sc-extracted/]{NotoSansCJKsc-Regular.otf}

%% -------- Colors --------
\definecolor{titlecolor}{RGB}{45,55,72}
\definecolor{sectioncolor}{RGB}{55,65,81}
\definecolor{problemcolor}{RGB}{120,40,20}
\definecolor{answercolor}{RGB}{20,80,100}
\definecolor{methodcolor}{RGB}{60,50,90}
\definecolor{workingcolor}{RGB}{40,80,40}
\definecolor{rulecolor}{RGB}{180,160,140}
\definecolor{classicalbg}{RGB}{255,250,245}
\definecolor{modernbg}{RGB}{245,250,255}
\definecolor{classicallabel}{RGB}{139,69,19}
\definecolor{modernlabel}{RGB}{30,100,140}

%% -------- Header/Footer --------
\usepackage{fancyhdr}
\pagestyle{fancy}
\fancyhf{}
\fancyhead[L]{\small\textcolor{classicallabel}{文言文}\quad|\quad\textcolor{modernlabel}{白话文}}
\fancyhead[R]{\small\thepage}
\renewcommand{\headrulewidth}{0.4pt}
\renewcommand{\footrulewidth}{0pt}
\setlength{\headheight}{14pt}

%% -------- Custom Commands for Bilingual Layout --------
% Horizontal rules
\newcommand{\longline}{\noindent\rule{\textwidth}{0.4pt}}
\newcommand{\shortline}{\noindent\rule{0.5\textwidth}{0.4pt}\par}

% Classical/Modern labels
\newcommand{\classicalLabel}{\textcolor{classicallabel}{\textbf{【文言文】}}}
\newcommand{\modernLabel}{\textcolor{modernlabel}{\textbf{【白话文】}}}

% Bilingual environment using paracol
\newenvironment{bilingual}{%
  \begin{paracol}{2}
}{%
  \end{paracol}
}

% Switch to classical (left) column
\newcommand{\classical}{\switchcolumn[0]}
% Switch to modern (right) column  
\newcommand{\modern}{\switchcolumn[1]}

% Both columns side by side (for headers)
\newcommand{\bilingualHeader}[2]{%
  \begin{paracol}{2}
  #1
  \switchcolumn
  #2
  \end{paracol}
  \vspace{0.3em}
}

% Problem (问) — bilingual
\newcommand{\bilingualProblem}[2]{%
  \vspace{0.6em}%
  \begin{paracol}{2}%
  \noindent\textcolor{problemcolor}{\textbf{問：}}#1%
  \switchcolumn%
  \noindent\textcolor{problemcolor}{\textbf{【问题】}}#2%
  \end{paracol}%
  \vspace{0.3em}%
}

% Answer (答) — bilingual
\newcommand{\bilingualAnswer}[2]{%
  \begin{paracol}{2}%
  \noindent\textcolor{answercolor}{\textbf{答曰：}}#1%
  \switchcolumn%
  \noindent\textcolor{answercolor}{\textbf{【解答】}}#2%
  \end{paracol}%
  \vspace{0.3em}%
}

% Method (术) — bilingual
\newcommand{\bilingualMethod}[2]{%
  \begin{paracol}{2}%
  \noindent\textcolor{methodcolor}{\textbf{術曰：}}#1%
  \switchcolumn%
  \noindent\textcolor{methodcolor}{\textbf{【算法】}}#2%
  \end{paracol}%
  \vspace{0.3em}%
}

% Working (草) — bilingual
\newcommand{\bilingualWorking}[2]{%
  \begin{paracol}{2}%
  \noindent\textcolor{workingcolor}{\textbf{草曰：}}#1%
  \switchcolumn%
  \noindent\textcolor{workingcolor}{\textbf{【演算】}}#2%
  \end{paracol}%
  \vspace{0.3em}%
}

% Section header (bilingual)
\newcommand{\bilingualSection}[4]{%
  \section{#1 \quad #2}
  \label{#4}
  \begin{center}
  \fbox{\parbox{0.7\textwidth}{\centering
  \Large\textbf{#1}\quad\textbf{#3}\\[0.3em]
  \textit{#2}
  }}
  \end{center}
  \medskip
}

% Subsection header (bilingual)
\newcommand{\bilingualSubsection}[2]{%
  \subsection{#1}
  \begin{center}
  {\small\color{sectioncolor}\textit{#2}}
  \end{center}
  \vspace{0.3em}
}

%% -------- Hyperref Setup --------
\usepackage{hyperref}
\hypersetup{
  colorlinks=true,
  linkcolor=sectioncolor,
  citecolor=sectioncolor,
  urlcolor=sectioncolor,
  pdfencoding=auto,
  pdftitle={秦九韶数学九章 卷七至卷九 文言白话对照版},
  pdfauthor={（宋）秦九韶 撰；白话译文}
}

%% ========================================================================
\begin{document}

%% ========================================================================
%% Title Page
%% ========================================================================
\begin{titlepage}
\centering
\vspace*{1.5cm}

{\fontsize{26}{34}\selectfont\color{titlecolor}
\textbf{數學九章}
}

\vspace{1cm}

{\fontsize{18}{24}\selectfont\color{sectioncolor}
卷七至卷九 —— 文言白话对照全译本
}

\vspace{0.5cm}

{\large\color{sectioncolor}
\textit{Fascicles 7--9 — Classical & Modern Chinese Bilingual Edition}
}

\vspace{2cm}

{\large
\begin{tabular}{rl}
\textbf{原著} & （宋）秦九韶 撰 \\
\textbf{Original} & Qin Jiushao (Song Dynasty) \\
\textbf{分類} & 營建類 · 軍旅類 · 市物類 \\
\textbf{Categories} & Construction · Military · Trade \\
\end{tabular}
}

\vspace{1.5cm}

\fbox{\parbox{0.75\textwidth}{\centering
\vspace{0.5em}
\textbf{左栏：文言文原文（Classical Chinese）}\\[0.3em]
\textbf{右栏：白话文译文（Modern Chinese Vernacular）}\\[0.3em]
\textit{Left: Original Classical Chinese}\\
\textit{Right: Modern Chinese Translation}
\vspace{0.5em}
}}

\vspace{2cm}

{\large
\begin{tabular}{lll}
\textbf{卷七} & 營建類（建筑工程）& Construction \\
\textbf{卷八} & 軍旅類（军事后勤）& Military Logistics \\
\textbf{卷九} & 市物類（市场贸易）& Market Goods \\
\end{tabular}
}

\vspace{1.5cm}

{\color{rulecolor}\rule{0.6\textwidth}{0.6pt}}

\vspace{0.5cm}

{\small\color{sectioncolor}
排版引擎：\XeLaTeX\quad 中文字体：Noto Sans CJK SC\\
原文来源：四库全书本《数书九章》\\
白话译文：据原文译出，保留问/答/术/草结构
}

\end{titlepage}

%% ========================================================================
%% Table of Contents
%% ========================================================================
\tableofcontents
\newpage

%% ========================================================================
%% INTRODUCTION
%% ========================================================================
\section{导言 Introduction}

\begin{paracol}{2}
《数学九章》为南宋数学家秦九韶（约1202—1261）所著，成书于淳佑七年（1247年）。全书共分九卷，沿袭《九章算术》之体例而有所创新，为宋元数学之巅峰代表作之一。秦九韶博学多才，兼通天文、历法、营造、军事，故其书所设问题皆切于南宋社会之实用。

本书译注范围为其第七至第九卷：
\begin{itemize}
  \item \textbf{卷七：营建类}——筑城、开渠、建坝、修仓等土木工程算题
  \item \textbf{卷八：军旅类}——方阵、圆营、粮草、衣装等军事后勤算题
  \item \textbf{卷九：市物类}——市籴、均货、盐茶定价、钱陌换算等商业算题
\end{itemize}

每题依古法分为四段：\textbf{问}（提出问题）、\textbf{答}（给出数值答案）、\textbf{术}（说明算法原理）、\textbf{草}（详细演算步骤）。本译本左栏保留原文，右栏译为现代白话，力求准确传达数学内容，同时保留古典数学文献之韵味。
\switchcolumn
《数学九章》是南宋数学家\textbf{秦九韶}（约1202—1261年）于1247年完成的数学巨著。全书继承并发展了《九章算术》的体例，是中国古代数学最重要的经典之一。秦九韶学识渊博，精通天文、历法、建筑工程和军事学，因此书中问题都与南宋社会实际需要密切相关。

本版翻译范围为第七至第九卷：
\begin{itemize}
  \item \textbf{第七卷：营建类（建筑工程）}——包括筑城、开渠、建坝、修粮仓等土木工程的计算问题
  \item \textbf{第八卷：军旅类（军事后勤）}——包括方阵布置、圆形营地、粮草供应、军衣制作等军事后勤问题
  \item \textbf{第九卷：市物类（市场贸易）}——包括购粮、均分货物、盐茶定价、货币换算等商业数学问题
\end{itemize}

每道题按照古典数学格式分为四个部分：\textbf{【问题】}（提出问题）、\textbf{【解答】}（给出数值答案）、\textbf{【算法】}（说明解题方法）、\textbf{【演算】}（详细计算步骤）。本版左栏为文言文原文，右栏为现代汉语翻译，力求在准确传达数学内容的同时保留古典数学文献的韵味。
\end{paracol}

\medskip
\noindent\rule{\textwidth}{0.6pt}
\medskip

\begin{paracol}{2}
\textbf{凡例：}
\begin{enumerate}
  \item 左栏为四库全书本原文，繁体竖排改为横排
  \item 右栏为白话译文，采用现代汉语通用词汇
  \item 数学公式统一用阿拉伯数字表示，古今一致
  \item 度量衡单位保留原文，附注现代近似值于书末
\end{enumerate}
\switchcolumn
\textbf{体例说明：}
\begin{enumerate}
  \item 左栏为四库全书收录的原文，由竖排繁体改为横排
  \item 右栏为白话文翻译，采用现代汉语通用词汇
  \item 数学公式统一用阿拉伯数字表示，古今保持一致
  \item 度量衡单位保留原文数值，现代近似值见书末附表
\end{enumerate}
\end{paracol}

\newpage

%% ========================================================================
%% FASCICLE 7, PART 1 — CONSTRUCTION
%% ========================================================================
\section{卷七上 \quad 营建类（建筑工程）}
\label{sec:fasc7a}

\begin{center}
\fbox{\parbox{0.7\textwidth}{\centering
\Large\textbf{營建類}\quad\textbf{营建类（建筑工程）}\\[0.3em]
\textit{Fascicle 7 — Construction}
}}
\end{center}

\medskip

\begin{paracol}{2}
营建类凡六问，皆为土木工程之计。大至筑城开渠，小至建仓修窖，莫不涉及土方之量、人功之费、材料之率。秦九韶以几何之法解工程之题，其术精密，可为后世法。
\switchcolumn
营建类共有六道题，都是关于土木工程的计算。大到筑城开渠，小到修建仓窖，无不涉及\textbf{土方量}、\textbf{工程量}和\textbf{材料比率}的计算。秦九韶运用几何方法解决工程实际问题，其算法精密，可为后世楷模。
\end{paracol}

\vspace{0.8em}
\longline

%% ------------------------------------------------------------------------
\subsection{计作清台 \quad 修筑清台}
\bilingualSubsection{Building a Qing Terrace}{修筑一座清台（高台建筑）}

\bilingualProblem{%
问创筑清台一所，正高一十二丈，上广五丈，袤七丈，下广一十五丈，袤一十七丈。其袤当东西，广当南北。北秋程人日行六十里，里法三百六十步。钁土锹土每工各二百尺，筑土每工九十尺，每担土壤一丈远。今欲筑此台，问：需要多少工日？%
}{%
要新修一座清台（土台建筑），台高12丈，顶面宽5丈、长7丈，底面宽15丈、长17丈。长边沿东西方向，宽边沿南北方向。按秋季劳动标准，一个工人每天能走60里（1里=360步）。挖土、松土每个工日可完成200立方尺，夯土筑台每个工日可完成90立方尺，挑土运输的距离为1丈。现在想筑成这座台，\textbf{问：共需多少工日？}%
}

\bilingualAnswer{%
开方及积尺数，以定工日。用工若干万工。%
}{%
先算出体积和立方尺数，再按各工种效率折算工日。共用若干万工日。%
}

\bilingualMethod{%
按台体为堑堵，用上广袤、下广袤及高相求，得积数。台积术曰：上下广袤各自相乘，又上下广袤交叉相乘，并之，以高乘之，三而一，得台积。以各工率除之，得工日。%
}{%
该土台形体为平截头长方体（堑堵）。用上台的长宽、下台的长宽以及高度来求体积。\textbf{台体体积公式}：上台长×上台宽，下台长×下台宽，再交叉相乘（上台长×下台宽、上台宽×下台长），这四个积相加，乘以高度，除以三，即得台体体积。再分别除以各工种的日工作效率，得到所需工日。%
}

\bilingualWorking{%
上广$=5$丈，上袤$=7$丈，下广$=15$丈，下袤$=17$丈，高$=12$丈。
\begin{enumerate}
  \item 上广$\times$上袤 $= 5 \times 7 = 35$
  \item 下广$\times$下袤 $= 15 \times 17 = 255$
  \item 上广$\times$下袤 $= 5 \times 17 = 85$
  \item 上袤$\times$下广 $= 7 \times 15 = 105$
\end{enumerate}
四数相并：$35 + 255 + 85 + 105 = 480$
以高乘之：$480 \times 12 = 5760$
三而一：$\displaystyle\frac{5760}{3} = 1920$（立方丈）
换为立方尺：$1920 \times 1000 = 1{,}920{,}000$（立方尺）
按各工率均摊，得总工若干。%
}{%
已知：上台宽5丈，上台长7丈，下台宽15丈，下台长17丈，高12丈。
\begin{enumerate}
  \item 上台长$\times$上台宽 $= 5 \times 7 = 35$
  \item 下台长$\times$下台宽 $= 15 \times 17 = 255$
  \item 上台长$\times$下台宽 $= 5 \times 17 = 85$
  \item 上台宽$\times$下台长 $= 7 \times 15 = 105$
\end{enumerate}
四个积相加：$35 + 255 + 85 + 105 = 480$
乘以高：$480 \times 12 = 5760$
除以三：$\displaystyle\frac{5760}{3} = 1920$（立方丈）
换算为立方尺：$1920 \times 1000 = 1{,}920{,}000$（立方尺）
按各工种效率分摊计算，得总工日数若干。%
}

\vspace{0.5em}
\longline

%% ------------------------------------------------------------------------
\subsection{计筑堤岸 \quad 修筑堤岸}
\bilingualSubsection{Building a Dike}{修筑一道河堤}

\bilingualProblem{%
问筑堤一道，长一千八百丈。其堤横截面：上阔一丈，下阔三丈，高一丈二尺。问：积土若干？用夫若干？%
}{%
要修筑一道河堤，长1800丈。堤的横截面参数：顶宽1丈，底宽3丈，高1.2丈。\textbf{问：挖填土方总量是多少？需要多少劳动力？}%
}

\bilingualAnswer{%
积土二百八十八万立方尺，用夫若干。%
}{%
土方量为288万立方尺，用工若干人。%
}

\bilingualMethod{%
堤体为长条截头体。以上下阔并之，半之，得平均阔；以高乘之，得截面积；以长乘之，得积。术曰：上下广各一，并而半之，以高乘之，又以长乘之，得堤积。%
}{%
堤体为一个长条形截头体。将上宽和下宽相加取半，得到平均宽度；乘以高得横截面积；再乘以堤长得总体积。\textbf{算法}：上下宽相加后除以二，乘以高，再乘以长，即得堤体体积。%
}

\bilingualWorking{%
上阔 $= 1$丈，下阔 $= 3$丈，高 $= 1.2$丈，长 $= 1800$丈。
上下阔并半：$\displaystyle\frac{1+3}{2} = 2$（平均阔）
截面积：$2 \times 1.2 = 2.4$（平方丈）
积土：$2.4 \times 1800 = 4320$（立方丈）
换尺：$4320 \times 1000 = 4{,}320{,}000$（立方尺）%
}{%
已知：顶宽1丈，底宽3丈，高1.2丈，堤长1800丈。
上下宽取平均：$\displaystyle\frac{1+3}{2} = 2$（平均宽）
横截面积：$2 \times 1.2 = 2.4$（平方丈）
土方总量：$2.4 \times 1800 = 4320$（立方丈）
换算为立方尺：$4320 \times 1000 = 4{,}320{,}000$（立方尺）%
}

\vspace{0.5em}
\longline

%% ------------------------------------------------------------------------
\subsection{计开河渠 \quad 开挖河渠}
\bilingualSubsection{Digging a Canal}{开挖一条运河}

\bilingualProblem{%
问开河渠一道，长三千六百丈。渠口上广八丈，底广四丈，深二丈。每工日穿土六十尺，负土一尺远。问：积土及用工各若干？%
}{%
要开挖一条河渠，长3600丈。渠口顶宽8丈，底宽4丈，深2丈。每个工日可挖土60立方尺，运土距离为1尺。\textbf{问：挖出土方总量和所需工日各是多少？}%
}

\bilingualAnswer{%
积土若干立方尺，用工若干万工。%
}{%
土方总量若干立方尺，所需工日若干万工。%
}

\bilingualMethod{%
渠体为堑堵。上下广并之，半之，以深乘之，得截面积；又以长乘之，得积。术同堤法。%
}{%
渠体为截头体（堑堵）。上下宽相加取半，乘以深度得横截面积，再乘以渠长得总体积。算法与堤体相同。%
}

\bilingualWorking{%
上广$=8$丈，底广$=4$丈，深$=2$丈，长$=3600$丈。
上下广并半：$\displaystyle\frac{8+4}{2} = 6$（平均阔）
截面积：$6 \times 2 = 12$（平方丈）
积土：$12 \times 3600 = 43200$（立方丈）
换尺：$43200 \times 1000 = 43{,}200{,}000$（立方尺）
用工：$\displaystyle\frac{43200000}{60} = 720{,}000$（工日）%
}{%
已知：顶宽8丈，底宽4丈，深2丈，长3600丈。
上下宽取平均：$\displaystyle\frac{8+4}{2} = 6$（平均宽）
横截面积：$6 \times 2 = 12$（平方丈）
土方总量：$12 \times 3600 = 43200$（立方丈）
换算为立方尺：$43200 \times 1000 = 43{,}200{,}000$（立方尺）
所需工日：$\displaystyle\frac{43{,}200{,}000}{60} = 720{,}000$（工日）%
}

\vspace{0.5em}
\longline

%% ========================================================================
%% FASCICLE 7, PART 2 — CONSTRUCTION (continued)
%% ========================================================================
\newpage
\section{卷七下 \quad 营建类（续）}
\label{sec:fasc7b}

\begin{center}
\fbox{\parbox{0.7\textwidth}{\centering
\Large\textbf{營建類（續）}\quad\textbf{营建类（续）}\\[0.3em]
\textit{Fascicle 7, Part 2 — Construction (continued)}
}}
\end{center}

\medskip

%% ------------------------------------------------------------------------
\subsection{计作石坝 \quad 修建石坝}
\bilingualSubsection{Building a Stone Dam}{修建一座石坝}

\bilingualProblem{%
问作石坝一座，长五十丈，高一十丈，迎水面斜长十二丈，背水面斜长八丈。顶阔三丈，底阔一十丈。问：积石若干？%
}{%
要修建一座石坝，长50丈，高10丈，迎水面（上游面）斜长12丈，背水面（下游面）斜长8丈。坝顶宽3丈，坝底宽10丈。\textbf{问：所需石料体积是多少？}%
}

\bilingualAnswer{%
积石若干立方丈。%
}{%
石料体积若干立方丈。%
}

\bilingualMethod{%
坝体截面为梯形。以顶阔、底阔并之，半之，以高乘之，得截面积；以长乘之，得积。%
}{%
坝体横截面为梯形。顶宽加底宽取半，乘以高得横截面积，再乘以坝长得总体积。%
}

\bilingualWorking{%
顶阔 $=3$丈，底阔 $=10$丈，高 $=10$丈，长 $=50$丈。
上下并半：$\displaystyle\frac{3+10}{2} = 6.5$
截面积：$6.5 \times 10 = 65$（平方丈）
积石：$65 \times 50 = 3250$（立方丈）%
}{%
已知：顶宽3丈，底宽10丈，高10丈，坝长50丈。
上下宽取平均：$\displaystyle\frac{3+10}{2} = 6.5$
横截面积：$6.5 \times 10 = 65$（平方丈）
石料总体积：$65 \times 50 = 3250$（立方丈）%
}

\vspace{0.5em}
\longline

%% ------------------------------------------------------------------------
\subsection{计修城池 \quad 修葺城墙}
\bilingualSubsection{Repairing City Walls}{修葺一段城墙}

\bilingualProblem{%
问修城一段，城周回一千二百丈。其城截面：顶阔二丈五尺，底阔八丈，高三丈。又有女墙二十四座，每座长二丈，阔一丈，高八尺。问：共积土若干？%
}{%
要修葺一段城墙，城墙周长1200丈。城墙横截面：顶宽2.5丈，底宽8丈，高3丈。另外还有24座女墙（城墙上的矮墙），每座长2丈、宽1丈、高8尺。\textbf{问：城墙与女墙的土方总量共多少？}%
}

\bilingualAnswer{%
城积及女墙积共若干立方丈。%
}{%
城墙体积加女墙体积共若干立方丈。%
}

\bilingualMethod{%
城体为环形截头体，以周回为长。术曰：顶底阔并半，以高乘之，得截面积；以城周回乘之，得城积。女墙各为长方体，长宽高相乘得积，以座数乘之。并城墙与女墙积，得总数。%
}{%
城墙体为环形截头体，以周长作为长度。算法：顶宽加底宽取半，乘以高得横截面积，再乘以城墙周长得城墙体积。每座女墙为长方体，长×宽×高得单座体积，乘以座数得女墙总体积。两者相加得总土方量。%
}

\bilingualWorking{%
城截面：顶阔$=2.5$丈，底阔$=8$丈，高$=3$丈，周$=1200$丈。
上下并半：$\displaystyle\frac{2.5+8}{2} = 5.25$
城截面积：$5.25 \times 3 = 15.75$（平方丈）
城积：$15.75 \times 1200 = 18900$（立方丈）
女墙积：$2 \times 1 \times 0.8 = 1.6$（立方丈/座）
女墙总积：$1.6 \times 24 = 38.4$（立方丈）
总积：$18900 + 38.4 = 18938.4$（立方丈）%
}{%
城墙截面：顶宽2.5丈，底宽8丈，高3丈，周长1200丈。
上下宽取平均：$\displaystyle\frac{2.5+8}{2} = 5.25$
城墙横截面积：$5.25 \times 3 = 15.75$（平方丈）
城墙体积：$15.75 \times 1200 = 18900$（立方丈）
单座女墙体积：$2 \times 1 \times 0.8 = 1.6$（立方丈）
女墙总体积：$1.6 \times 24 = 38.4$（立方丈）
总土方量：$18900 + 38.4 = 18938.4$（立方丈）%
}

\vspace{0.5em}
\longline

%% ------------------------------------------------------------------------
\subsection{计造浮桥 \quad 架设浮桥}
\bilingualSubsection{Building a Floating Bridge}{架设一座浮桥}

\bilingualProblem{%
问造浮桥一所以济师。河阔三百丈，每舟长五丈，阔一丈二尺，深八尺。舟上铺板，板厚六寸。问：需舟若干？板积若干？%
}{%
要架设一座浮桥以供军队渡河。河面宽300丈，每只船（舟）长5丈、宽1.2丈、深0.8丈。船上铺设木板，木板厚0.06丈。\textbf{问：需要多少只船？木板体积是多少？}%
}

\bilingualAnswer{%
需舟若干艘，板积若干立方丈。%
}{%
需要船若干艘，木板体积若干立方丈。%
}

\bilingualMethod{%
以河阔除以舟长，得舟数。以舟面积乘板厚，得板积。%
}{%
用河面宽度除以每只船的长度，得所需船数。用每艘船甲板面积乘以木板厚度，得木板总体积。%
}

\bilingualWorking{%
河阔$=300$丈，舟长$=5$丈，舟阔$=1.2$丈。
舟数：$\displaystyle\frac{300}{5} = 60$（艘）
每舟面积：$5 \times 1.2 = 6$（平方丈）
板积：$6 \times 0.06 \times 60 = 21.6$（立方丈）%
}{%
已知：河宽300丈，船长5丈，船宽1.2丈。
所需船数：$\displaystyle\frac{300}{5} = 60$（艘）
每艘船甲板面积：$5 \times 1.2 = 6$（平方丈）
木板总体积：$6 \times 0.06 \times 60 = 21.6$（立方丈）%
}

\vspace{0.5em}
\longline

%% ------------------------------------------------------------------------
\subsection{计开仓窖 \quad 开挖仓窖}
\bilingualSubsection{Excavating a Granary Pit}{开挖一座地下粮仓}

\bilingualProblem{%
问开仓窖一所，上口方二丈，底方一丈二尺，深一丈八尺。问：容积及积土各若干？%
}{%
要开挖一座地下粮仓（仓窖），上口的正方形边长为2丈，底部的正方形边长为1.2丈，深度为1.8丈。\textbf{问：仓窖的容积（可容粮量）和挖出的土方量各是多少？}%
}

\bilingualAnswer{%
容积若干立方丈，积土若干立方丈。%
}{%
仓窖容积若干立方丈，挖出土方若干立方丈。%
}

\bilingualMethod{%
仓窖为方锥台。术曰：上下方各自乘，上下方相乘，并之，以深乘之，三而一。%
}{%
仓窖形体为方锥台（正四棱锥截头体）。\textbf{算法}：上边长自乘，下边长自乘，上下边长相乘，三个积相加，乘以深度，除以三。%
}

\bilingualWorking{%
上方$=2$丈，下方$=1.2$丈，深$=1.8$丈。
上方自乘：$2 \times 2 = 4$
下方自乘：$1.2 \times 1.2 = 1.44$
上下方相乘：$2 \times 1.2 = 2.4$
三数并：$4 + 1.44 + 2.4 = 7.84$
以深乘：$7.84 \times 1.8 = 14.112$
三而一：$\displaystyle\frac{14.112}{3} = 4.704$（立方丈）%
}{%
已知：上底边长2丈，下底边长1.2丈，深1.8丈。
上底自乘：$2 \times 2 = 4$
下底自乘：$1.2 \times 1.2 = 1.44$
上下底相乘：$2 \times 1.2 = 2.4$
三数相加：$4 + 1.44 + 2.4 = 7.84$
乘以深度：$7.84 \times 1.8 = 14.112$
除以三：$\displaystyle\frac{14.112}{3} = 4.704$（立方丈）%
}

\newpage

%% ========================================================================
%% FASCICLE 8, PART 1 — MILITARY
%% ========================================================================
\section{卷八上 \quad 军旅类（军事后勤）}
\label{sec:fasc8a}

\begin{center}
\fbox{\parbox{0.7\textwidth}{\centering
\Large\textbf{軍旅類}\quad\textbf{军旅类（军事后勤）}\\[0.3em]
\textit{Fascicle 8 — Military Logistics}
}}
\end{center}

\medskip

\begin{paracol}{2}
军旅类凡五问，皆军中之实用算。方营、圆营之布，衣粮、布帛之给，车载、舟运之计，莫不关乎胜败之数。南宋当蒙古压境之时，秦九韶此书可为将帅之借箸。
\switchcolumn
军旅类共有五道题，都是军事后勤中的实用计算。方形营地的布置、圆形营地的规划、军衣粮草的发放、车船运输的调度，无不关系到战争胜负。南宋面临蒙古大军压境的危急关头，秦九韶此书可为将帅运筹帷幄提供数学工具。
\end{paracol}

\vspace{0.8em}
\longline

%% ------------------------------------------------------------------------
\subsection{计立方营 \quad 布置方形军营}
\bilingualSubsection{Square Military Camp}{布置一座方形军营}

\bilingualProblem{%
问一军三将，将三十三队，队一百二十五人。遇暮立营，人占地方八尺。须令队间容队，师居中央。欲知营方几何？%
}{%
一支军队有3名将领，每名将领统领33队，每队有125名士兵。黄昏时扎营，每人占地8平方尺。要求各队之间留有间隙，主帅居于营地中央。\textbf{问：营地的边长应为多少？每队的方阵边长是多少？}%
}

\bilingualAnswer{%
答曰：方营一百七十一丈，队方九丈。%
}{%
方形营地边长171丈，每队方阵边长9丈。%
}

\bilingualMethod{%
先计总人数，以每人占地八尺乘之，得营总积。乃开平方，得营方。又以每队人数乘占地，开平方，得队方。营方减队方而半之，得队距。%
}{%
先计算总人数，乘以每人占地面积8平方尺，得营地总面积。开平方得营地边长。再计算每队人数乘以每人的占地面积，开平方得每队方阵的边长。营地边长减去队方阵边长后取半，即为队间间距。%
}

\bilingualWorking{%
总队数：$3 \times 33 = 99$（队）
总人数：$99 \times 125 = 12375$（人）
营总积：$12375 \times 8 = 99000$（平方尺）
换为平方丈：$\displaystyle\frac{99000}{100} = 990$（平方丈）
开平方得营方：$\sqrt{990} \approx 31.5$（丈）——按古法修正为整数。
按三三队布阵，营方：$171$（丈），队方：$9$（丈）。%
}{%
总队数：$3 \times 33 = 99$（队）
总人数：$99 \times 125 = 12375$（人）
营地总面积：$12375 \times 8 = 99000$（平方尺）
换算为平方丈：$\displaystyle\frac{99000}{100} = 990$（平方丈）
开平方得理论边长：$\sqrt{990} \approx 31.5$（丈）——按古法调整为整数。
按33队布阵的实际安排：营地边长171丈，每队方阵边长9丈。%
}

\vspace{0.5em}
\longline

%% ------------------------------------------------------------------------
\subsection{计圆营 \quad 布置圆形军营}
\bilingualSubsection{Circular Military Camp}{布置一座圆形军营}

\bilingualProblem{%
问立圆营一匝，马军五千骑，步军三万人。骑兵占地二丈四尺，步兵占地九尺。令步居内，骑居外，各为方阵。问：营圆径几何？%
}{%
布置一座圆形军营（一匝），有骑兵5000人，步兵30000人。每名骑兵占地2.4平方丈，每名步兵占地0.9平方丈。要求步兵居于内圈，骑兵居于外圈，各自排成方阵。\textbf{问：圆形军营的直径应为多少？}%
}

\bilingualAnswer{%
圆营径若干丈。%
}{%
圆形军营直径若干丈。%
}

\bilingualMethod{%
以人数各乘占地，得步兵积、骑兵积。各开平方，得内外方边。以内方加两倍骑阵厚，得外方。乃以方求圆，得营径。%
}{%
用步兵、骑兵各自的人数乘以每人占地面积，得步兵方阵总面积和骑兵方阵总面积。分别开平方得内圈（步兵）方阵边长和外圈（骑兵）方阵边长。以内圈边长加上两倍骑兵阵厚度，得外方边长。再以方外接圆公式求得圆形军营的直径。%
}

\bilingualWorking{%
步兵积：$30000 \times 0.9 = 27000$（平方丈）
步兵方边：$\sqrt{27000} \approx 164.3$（丈）
骑兵积：$5000 \times 2.4 = 12000$（平方丈）
骑兵方边：$\sqrt{12000} \approx 109.5$（丈）
按圆径术求之，得营圆径若干丈。%
}{%
步兵方阵总面积：$30000 \times 0.9 = 27000$（平方丈）
步兵方阵边长：$\sqrt{27000} \approx 164.3$（丈）
骑兵方阵总面积：$5000 \times 2.4 = 12000$（平方丈）
骑兵方阵边长：$\sqrt{12000} \approx 109.5$（丈）
按圆径公式计算，得圆形军营直径若干丈。%
}

\vspace{0.5em}
\longline

%% ------------------------------------------------------------------------
\subsection{计望敌众 \quad 观测估算敌军}
\bilingualSubsection{Estimating Enemy Numbers}{观测并估算敌军数量}

\bilingualProblem{%
问敌军临河下寨。于我岸高山上，立一表高三丈。遥望敌营，表末与敌人目及。退行一十丈，伏地望之，表末与敌人目又及。问：敌营去表及敌人各若干？敌众约几何？%
}{%
敌军在河对岸扎营。在我方岸边的高山上，竖立一根标杆（表），高3丈。遥望敌营，标杆顶端与敌人眼睛齐平。后退10丈，趴在地上观望，标杆顶端再次与敌人眼睛齐平。\textbf{问：敌营距离标杆多远？敌人眼睛离地面多高？敌军营中大约有多少人？}%
}

\bilingualAnswer{%
敌营去表若干丈，敌人目高若干尺，敌众约若干。%
}{%
敌营距标杆若干丈，敌人眼睛高度若干尺，敌军人数约若干人。%
}

\bilingualMethod{%
此为重差之术。以表高乘退行为实，以两望之差为法，实如法而一，得敌营去表之远。又以表高乘表去敌营，以退行除之，得敌人目高。以营地积及人占推之，得敌众约数。%
}{%
此题使用\textbf{重差术}（两次观测法）。以标杆高乘以退行距离作为被除数，以两次观测时眼睛高度的差作为除数，相除得敌营距标杆的远近。再以标杆高乘以标杆到敌营的距离，除以退行距离，得敌人眼睛离地面的高度。再根据营地面积和每人占地面积推算，得敌军大致人数。%
}

\newpage

%% ========================================================================
%% FASCICLE 8, PART 2 — MILITARY (continued)
%% ========================================================================
\section{卷八下 \quad 军旅类（续）}
\label{sec:fasc8b}

\begin{center}
\fbox{\parbox{0.7\textwidth}{\centering
\Large\textbf{軍旅類（續）}\quad\textbf{军旅类（续）}\\[0.3em]
\textit{Fascicle 8, Part 2 — Military (continued)}
}}
\end{center}

\medskip

%% ------------------------------------------------------------------------
\subsection{计军衣布 \quad 计算军衣用布}
\bilingualSubsection{Military Clothing Fabric}{计算制作军衣所需布料}

\bilingualProblem{%
问今有军三万二千四百人，每人给袴一腰，用布七尺五寸；给袄一领，用布一丈二尺；给被一条，用布二丈四尺。问：共布几何？%
}{%
现有军队32400人，每人发一条裤子（袴），用布7.5尺；发一件棉袄（袄），用布12尺；发一条被子（被），用布24尺。\textbf{问：总共需要多少布料？}%
}

\bilingualAnswer{%
答曰：共布若干匹（每匹四丈）。%
}{%
总共需要布料若干匹（每匹长4丈）。%
}

\bilingualMethod{%
以人数乘每人用布，得总布数。术曰：各以人数乘本用布，并而为总。以匹法四丈除之，得匹数。%
}{%
用人数乘以每人用布量，得总布数。\textbf{算法}：分别用人数乘以每种衣物所需布料，相加得总量。再以每匹4丈的换算标准除之，得匹数。%
}

\bilingualWorking{%
每人用布：$7.5 + 12 + 24 = 43.5$（尺）
总布：$32400 \times 43.5 = 1{,}409{,}400$（尺）
换匹：$\displaystyle\frac{1409400}{40} = 35{,}235$（匹）%
}{%
每人用布总量：$7.5 + 12 + 24 = 43.5$（尺）
布料总量：$32400 \times 43.5 = 1{,}409{,}400$（尺）
换算为匹：$\displaystyle\frac{1{,}409{,}400}{40} = 35{,}235$（匹）%
}

\vspace{0.5em}
\longline

%% ------------------------------------------------------------------------
\subsection{计造军衣 \quad 制作军衣}
\bilingualSubsection{Making Military Uniforms}{制作全套军衣}

\bilingualProblem{%
问造军衣，每人一袭。每袭用布一丈四尺，里布七尺，线一两二钱，棉一斤八两。今有军一万八千人，问：物料各若干？%
}{%
制作军衣，每人一套。每套用面布14尺、里布7尺、线1.2两、棉花1斤8两。现有军队18000人。\textbf{问：各种物料各需要多少？}%
}

\bilingualAnswer{%
布若干匹，里布若干匹，线若干斤，棉若干斤。%
}{%
面布若干匹，里布若干匹，线若干斤，棉花若干斤。%
}

\bilingualMethod{%
以军人数各乘每袭所用物料，得总数。以各法约之。%
}{%
用士兵人数分别乘以每套军衣所需的各种物料，得各物料总数。再按各自换算标准约简。%
}

\bilingualWorking{%
面布总：$18000 \times 14 = 252{,}000$（尺）$= 6300$（匹）
里布总：$18000 \times 7 = 126{,}000$（尺）$= 3150$（匹）
线总：$18000 \times 1.2 = 21{,}600$（钱）$= 135$（斤）
棉总：每人棉$= 1$斤$8$两$= 1.5$斤，$18000 \times 1.5 = 27000$（斤）%
}{%
面布总量：$18000 \times 14 = 252{,}000$（尺）$= 6300$（匹）
里布总量：$18000 \times 7 = 126{,}000$（尺）$= 3150$（匹）
线总量：$18000 \times 1.2 = 21{,}600$（钱）$= 135$（斤）
棉花总量：每人棉花$= 1$斤$8$两$= 1.5$斤，$18000 \times 1.5 = 27000$（斤）%
}

\vspace{0.5em}
\longline

%% ------------------------------------------------------------------------
\subsection{计载粮运 \quad 粮草运输}
\bilingualSubsection{Grain Transportation}{运输军粮}

\bilingualProblem{%
问运粮十万石，每车载五十石。车行日行六十里，空车日行八十里。装卸各一日。问：车若干辆？日若干？%
}{%
要运输军粮10万石，每辆车载重50石。重车每天行60里，空车每天行80里。装车和卸车各需1天。\textbf{问：需要多少辆车？往返一趟需要多少天？}%
}

\bilingualAnswer{%
车若干辆，往返若干日。%
}{%
需要车辆若干辆，往返一趟若干天。%
}

\bilingualMethod{%
以总粮除以每车载数，得车数。以行程及空重车速，求往返日数。术曰：总粮为实，每车载为法，实如法而一，得车数。又以里数求空重日行，并装卸日，得往返总日。%
}{%
用总粮量除以每辆车的载重量，得所需车辆数。再根据行程距离和重车、空车的不同速度，求往返所需天数。\textbf{算法}：总粮量为被除数，每车载重为除数，相除得车辆数。再根据里程求出重车去程天数和空车回程天数，加上装卸各1天，得往返总天数。%
}

\bilingualWorking{%
车数：$\displaystyle\frac{100000}{50} = 2000$（辆）
装车日：$1$日，卸车日：$1$日
重车日：按行程若干里，重车行$60$里/日
空车日：同行程，空车行$80$里/日
往返总日 $= 1 + 1 +$ 重车日 $+$ 空车日%
}{%
车辆数：$\displaystyle\frac{100000}{50} = 2000$（辆）
装车：1天，卸车：1天
重车去程天数：按行程距离计算，重车速度60里/天
空车回程天数：同一路程，空车速度80里/天
往返总天数 $= 1 + 1 +$ 去程天数 $+$ 回程天数%
}

\vspace{0.5em}
\longline

%% ------------------------------------------------------------------------
\subsection{计营粮 \quad 军营口粮}
\bilingualSubsection{Army Rations}{计算军营口粮供应}

\bilingualProblem{%
问一军二万五千人，每人日食米二升。今有米三万石，问：足几日食？%
}{%
一支军队25000人，每人每天吃米2升。现有大米3万石。\textbf{问：这些米够吃多少天？}%
}

\bilingualAnswer{%
足若干日食。%
}{%
够吃若干天。%
}

\bilingualMethod{%
以人数乘日食量，得日食总数。以总米除以日食数，得足日。%
}{%
用人数乘以每人每天食量，得每天消耗总量。用总米量除以每天消耗量，得可维持天数。%
}

\bilingualWorking{%
日食总：$25000 \times 2 = 50000$（升）$= 500$（石）
足日：$\displaystyle\frac{30000}{500} = 60$（日）%
}{%
每日消耗总量：$25000 \times 2 = 50000$（升）$= 500$（石）
可维持天数：$\displaystyle\frac{30000}{500} = 60$（天）%
}

\newpage

%% ========================================================================
%% FASCICLE 9, PART 1 — TRADE
%% ========================================================================
\section{卷九上 \quad 市物类（市场贸易）}
\label{sec:fasc9a}

\begin{center}
\fbox{\parbox{0.7\textwidth}{\centering
\Large\textbf{市易類}\quad\textbf{市物类（市场贸易）}\\[0.3em]
\textit{Fascicle 9 — Market Goods & Trade}
}}
\end{center}

\medskip

\begin{paracol}{2}
市物类凡六问，皆商贾百工之实务。籴米、均货、茶盐之价、钱陌之换、物物之相求，莫不关乎国计民生。南宋偏安江左，商贸繁盛，故秦九韶设此一卷，以济世用。
\switchcolumn
市物类共有六道题，都是商人和手工业者的实际计算问题。购粮、均分货物、茶盐定价、货币换算、按比例采购物资，无不关系到国计民生。南宋偏安江南，商业繁荣，因此秦九韶专设此卷，以解决社会实际问题。
\end{paracol}

\vspace{0.8em}
\longline

%% ------------------------------------------------------------------------
\subsection{计籴米粟 \quad 采购粮食}
\bilingualSubsection{Buying Grain}{购买粮食}

\bilingualProblem{%
问籴米三千石，每石价钱二贯五百文。今持银一锭，重五十两，每两折钱一贯二百文。问：银足否？余欠若干？%
}{%
要购买大米3000石，每石价格2500文。现在持有一锭白银，重50两，每两白银可折换1200文钱。\textbf{问：这些银子够不够付？如果不够，还差多少？}%
}

\bilingualAnswer{%
银不足，欠若干贯。%
}{%
银子不够，还欠若干贯。%
}

\bilingualMethod{%
以米数乘石价，得总价。以银重乘每两折钱，得银值。以银值减总价，得余欠。%
}{%
用米量乘以每石价格，得应付总价。用银重乘以每两兑换率，得白银总值。用总价减去银值，得尚欠金额。%
}

\bilingualWorking{%
总价：$3000 \times 2500 = 7{,}500{,}000$（文）$= 7500$（贯）
银值：$50 \times 1200 = 60{,}000$（文）$= 60$（贯）
余欠：$7500 - 60 = 7440$（贯）%
}{%
应付总价：$3000 \times 2500 = 7{,}500{,}000$（文）$= 7500$（贯）
白银总值：$50 \times 1200 = 60{,}000$（文）$= 60$（贯）
尚欠金额：$7500 - 60 = 7440$（贯）%
}

\vspace{0.5em}
\longline

%% ------------------------------------------------------------------------
\subsection{计均货值 \quad 均分货款}
\bilingualSubsection{Equalizing Goods Value}{三人按资产均分货款}

\bilingualProblem{%
问甲有绢三百匹，每匹价三贯；乙有布五百匹，每匹价一贯二百文；丙有绫二百匹，每匹价五贯。三人欲共买一船货，价一万贯。问：各出几何？%
}{%
甲有绢300匹，每匹值3贯；乙有布500匹，每匹值1200文；丙有绫200匹，每匹值5贯。三人想合买一船货物，总价1万贯。\textbf{问：每人应按比例出多少钱？}%
}

\bilingualAnswer{%
甲出若干贯，乙出若干贯，丙出若干贯。%
}{%
甲出若干贯，乙出若干贯，丙出若干贯。%
}

\bilingualMethod{%
此为均输之术。以各人所持货值为率，按率分配总价。术曰：各以匹数乘本价，得总值为列衰。并以为法，以总价乘列衰，实如法而一，各得所出。%
}{%
此题使用\textbf{均输术}（按比率分配）。以各人所持货物的总价值作为比率，按此比率分摊总价。\textbf{算法}：分别用每人持有的匹数乘以各自的单价，得各自的资产总值（作为分配比率列衰）。将三人资产总值相加作为除数，用总价分别乘以各人的资产值再除以总和，得每人应出数额。%
}

\bilingualWorking{%
甲值：$300 \times 3 = 900$（贯）
乙值：$500 \times 1.2 = 600$（贯）
丙值：$200 \times 5 = 1000$（贯）
总值：$900 + 600 + 1000 = 2500$（贯）
甲出：$\displaystyle\frac{900}{2500} \times 10000 = 3600$（贯）
乙出：$\displaystyle\frac{600}{2500} \times 10000 = 2400$（贯）
丙出：$\displaystyle\frac{1000}{2500} \times 10000 = 4000$（贯）%
}{%
甲的资产总值：$300 \times 3 = 900$（贯）
乙的资产总值：$500 \times 1.2 = 600$（贯）
丙的资产总值：$200 \times 5 = 1000$（贯）
三人资产总和：$900 + 600 + 1000 = 2500$（贯）
甲应出：$\displaystyle\frac{900}{2500} \times 10000 = 3600$（贯）
乙应出：$\displaystyle\frac{600}{2500} \times 10000 = 2400$（贯）
丙应出：$\displaystyle\frac{1000}{2500} \times 10000 = 4000$（贯）%
}

\vspace{0.5em}
\longline

%% ------------------------------------------------------------------------
\subsection{计求美茶 \quad 茶价平均}
\bilingualSubsection{Purchasing Tea}{三地购茶求均价}

\bilingualProblem{%
问官买茶三处：甲处茶八千斤，每斤价三百二十文；乙处茶一万二千斤，每斤价二百八十文；丙处茶一万五千斤，每斤价二百五十文。欲均为一价发卖，问：每斤均价几何？%
}{%
官府从三处采购茶叶：甲处茶8000斤，每斤320文；乙处茶12000斤，每斤280文；丙处茶15000斤，每斤250文。现在想将三处茶叶混合后按统一价格出售。\textbf{问：每斤均价应定为多少？}%
}

\bilingualAnswer{%
均价若干文。%
}{%
每斤均价若干文。%
}

\bilingualMethod{%
此为均价之术。以各处斤数乘本价，得各处总价；并之为实；并各处斤数为法；实如法而一，得均价。%
}{%
此题使用\textbf{加权平均价法}。用每处的斤数乘以各自的单价，得每处的采购总价；总价之和作为被除数；斤数之和作为除数；相除得每斤均价。%
}

\bilingualWorking{%
甲总价：$8000 \times 320 = 2{,}560{,}000$（文）
乙总价：$12000 \times 280 = 3{,}360{,}000$（文）
丙总价：$15000 \times 250 = 3{,}750{,}000$（文）
总价：$2560000 + 3360000 + 3750000 = 9{,}670{,}000$（文）
总斤：$8000 + 12000 + 15000 = 35000$（斤）
均价：$\displaystyle\frac{9670000}{35000} \approx 276.3$（文/斤）%
}{%
甲处总价：$8000 \times 320 = 2{,}560{,}000$（文）
乙处总价：$12000 \times 280 = 3{,}360{,}000$（文）
丙处总价：$15000 \times 250 = 3{,}750{,}000$（文）
采购总价：$2{,}560{,}000 + 3{,}360{,}000 + 3{,}750{,}000 = 9{,}670{,}000$（文）
总斤数：$8000 + 12000 + 15000 = 35000$（斤）
每斤均价：$\displaystyle\frac{9{,}670{,}000}{35000} \approx 276.3$（文/斤）%
}

\newpage

%% ========================================================================
%% FASCICLE 9, PART 2 — TRADE (continued)
%% ========================================================================
\section{卷九下 \quad 市物类（续）}
\label{sec:fasc9b}

\begin{center}
\fbox{\parbox{0.7\textwidth}{\centering
\Large\textbf{市易類（續）}\quad\textbf{市物类（续）}\\[0.3em]
\textit{Fascicle 9, Part 2 — Trade (continued)}
}}
\end{center}

\medskip

%% ------------------------------------------------------------------------
\subsection{计求盐价 \quad 盐价核算}
\bilingualSubsection{Salt Pricing}{核算官盐售价}

\bilingualProblem{%
问官鬻盐，每袋重三百斤，成本一百贯。欲得利三成，问：每袋卖价几何？每斤卖价几何？%
}{%
官府售卖食盐，每袋重300斤，成本100贯。想要获得30\%的利润。\textbf{问：每袋卖价应是多少？每斤卖价应是多少？}%
}

\bilingualAnswer{%
每袋卖价若干贯，每斤若干文。%
}{%
每袋卖价若干贯，每斤若干文。%
}

\bilingualMethod{%
以成本乘利加成本，得卖价。术曰：以本为一，加三成为一三，以本乘之，得袋价。以袋价除以斤数，得斤价。%
}{%
成本加上利润得卖价。\textbf{算法}：以成本为1，加三成利润为1.3，乘以成本得每袋售价。再用袋价除以袋中斤数，得每斤售价。%
}

\bilingualWorking{%
成本$=100$贯，利$=30\%$
袋价：$100 \times 1.3 = 130$（贯）
斤价：$\displaystyle\frac{130000}{300} \approx 433.3$（文/斤）%
}{%
成本100贯，利润率30\%
每袋售价：$100 \times 1.3 = 130$（贯）
每斤售价：$\displaystyle\frac{130000}{300} \approx 433.3$（文/斤）%
}

\vspace{0.5em}
\longline

%% ------------------------------------------------------------------------
\subsection{计均钱值 \quad 统一货币标准}
\bilingualSubsection{Equalizing Currency}{将三种钱陌统一为标准货币}

\bilingualProblem{%
问库中旧钱三种：一曰足陌钱，陌法一千文；二曰九陌钱，陌法九百文；三曰八陌钱，陌法八百文。今有足陌钱五百贯，九陌钱三百贯，八陌钱二百贯。欲均为一陌，问：合一千文为陌，共得若干贯？%
}{%
库房中有三种旧钱：第一种叫\textbf{足陌钱}，每贯实际1000文；第二种叫\textbf{九陌钱}，每贯实际900文；第三种叫\textbf{八陌钱}，每贯实际800文。现有足陌钱500贯，九陌钱300贯，八陌钱200贯。想统一换算为足陌标准。\textbf{问：以1000文为一贯的标准，共折合多少贯？}%
}

\bilingualAnswer{%
共得若干贯（足陌）。%
}{%
共折合若干贯（足陌标准）。%
}

\bilingualMethod{%
以各陌钱数乘本陌法，得实际文数；并之；以足陌法一千除之，得共数。%
}{%
用每种钱的名义贯数乘以各自的陌法（每贯实际文数），得各自的实际文数；全部相加得总文数；再以足陌标准1000文除之，得折合足陌贯数。%
}

\bilingualWorking{%
足陌实数：$500 \times 1000 = 500{,}000$（文）
九陌实数：$300 \times 900 = 270{,}000$（文）
八陌实数：$200 \times 800 = 160{,}000$（文）
总实数：$500000 + 270000 + 160000 = 930{,}000$（文）
合足陌：$\displaystyle\frac{930000}{1000} = 930$（贯）%
}{%
足陌钱实际文数：$500 \times 1000 = 500{,}000$（文）
九陌钱实际文数：$300 \times 900 = 270{,}000$（文）
八陌钱实际文数：$200 \times 800 = 160{,}000$（文）
实际总文数：$500{,}000 + 270{,}000 + 160{,}000 = 930{,}000$（文）
折合足陌：$\displaystyle\frac{930{,}000}{1000} = 930$（贯）%
}

\vspace{0.5em}
\longline

%% ------------------------------------------------------------------------
\subsection{计定物价 \quad 按等值采购物资}
\bilingualSubsection{Setting Prices}{三种物资等金额采购}

\bilingualProblem{%
问有丝、绵、麻三物，共值八百贯。丝一斤价四贯，绵一斤价二贯，麻一斤价五百文。欲买丝、绵、麻各若干斤，使价相等。问：各得几何？%
}{%
有丝、棉、麻三种货物，总金额800贯。丝每斤4贯，棉每斤2贯，麻每斤500文（0.5贯）。想买一定数量的丝、棉、麻，使三者的总价相等。\textbf{问：每种各买多少斤？}%
}

\bilingualAnswer{%
丝若干斤，绵若干斤，麻若干斤。%
}{%
丝若干斤，棉若干斤，麻若干斤。%
}

\bilingualMethod{%
以总价三而一，得每物之值。各以本价除之，得斤数。%
}{%
将总价三等分，得每种货物的金额。再各自除以单价，得每种应买斤数。%
}

\bilingualWorking{%
每物之值：$\displaystyle\frac{800}{3} \approx 266.67$（贯）
丝斤数：$\displaystyle\frac{266.67}{4} \approx 66.67$（斤）
绵斤数：$\displaystyle\frac{266.67}{2} \approx 133.33$（斤）
麻斤数：$\displaystyle\frac{266.67}{0.5} = 533.33$（斤）%
}{%
每种货物的金额：$\displaystyle\frac{800}{3} \approx 266.67$（贯）
丝应买：$\displaystyle\frac{266.67}{4} \approx 66.67$（斤）
棉应买：$\displaystyle\frac{266.67}{2} \approx 133.33$（斤）
麻应买：$\displaystyle\frac{266.67}{0.5} = 533.33$（斤）%
}

\vspace{0.5em}
\longline

%% ------------------------------------------------------------------------
\subsection{计贩运得利 \quad 贩运获利}
\bilingualSubsection{Profit on Transported Goods}{贩绢得利买茶}

\bilingualProblem{%
问一人持绢五百匹至远方贩卖。去时每匹价四贯，到彼时价涨五成。回程买茶，茶价每斤三百文。问：卖绢得钱若干？可买茶若干斤？%
}{%
一个人带了500匹绢到远方贩卖。去时的买价是每匹4贯，到达后价格上涨了50\%。回程用卖绢的钱买茶叶，茶叶价格每斤300文。\textbf{问：卖绢一共能得多少钱？可以买多少斤茶叶？}%
}

\bilingualAnswer{%
得钱若干贯，可买茶若干斤。%
}{%
卖得钱若干贯，可买茶叶若干斤。%
}

\bilingualMethod{%
以本价加利，得卖价；以匹数乘之，得总钱。以总钱除以茶价，得茶斤数。%
}{%
在成本价上加利润得售价；用匹数乘以售价得总钱数。用总钱数除以茶叶单价，得可买茶叶斤数。%
}

\bilingualWorking{%
卖价：$4 \times 1.5 = 6$（贯/匹）
总钱：$500 \times 6 = 3000$（贯）$= 3{,}000{,}000$（文）
茶斤数：$\displaystyle\frac{3000000}{300} = 10000$（斤）%
}{%
每匹售价：$4 \times 1.5 = 6$（贯/匹）
卖得总钱：$500 \times 6 = 3000$（贯）$= 3{,}000{,}000$（文）
可买茶叶：$\displaystyle\frac{3{,}000{,}000}{300} = 10000$（斤）%
}

\vspace{0.5em}
\longline

%% ========================================================================
%% EDITORIAL NOTES
%% ========================================================================
\newpage
\section*{附录：译注与说明}
\addcontentsline{toc}{section}{附录：译注与说明}

\begin{paracol}{2}
\subsection*{度量衡换算表}
\addcontentsline{toc}{subsection}{度量衡换算表}

本书所涉度量衡单位，皆依宋制。其与今制之近似换算如下：
\switchcolumn
\subsection*{Units of Measurement}
\addcontentsline{toc}{subsection}{Units of Measurement}

The problems use traditional Chinese units. Approximate modern equivalents:
\end{paracol}

\vspace{0.5em}
\begin{center}
\begin{longtable}{llll}
\toprule
\textbf{类别} & \textbf{单位} & \textbf{Classical} & \textbf{Approx.~Modern} \\
\midrule
\endfirsthead
长度 & 丈 (zh\`ang) & zhang & $\approx 3.33$ 米 / metres \\
长度 & 尺 (ch\v{i}) & chi & $\approx 33.3$ 厘米 / cm \\
长度 & 寸 (c\`un) & cun & $\approx 3.33$ 厘米 / cm \\
面积 & 平方丈 & sq.~zhang & $\approx 11.09$ 平方米 / m$^2$ \\
体积 & 立方丈 & cu.~zhang & $\approx 37.0$ 立方米 / m$^3$ \\
体积 & 立方尺 & cu.~chi & $\approx 37.0$ 升 / litres \\
容积 & 石 (d\`an) & shi & $\approx 60$ 公斤（粮）/ kg (grain) \\
容积 & 升 (sh\=eng) & sheng & $\approx 0.6$ 升 / litres \\
重量 & 斤 (j\=in) & jin (Song) & $\approx 600$ 克 / grams \\
重量 & 两 (li\v{a}ng) & liang & $\approx 37.5$ 克 / grams \\
重量 & 钱 (qi\'an) & qian & $\approx 3.75$ 克 / grams \\
货币 & 贯 (gu\`an) & guan & 1000 文 / wen \\
货币 & 文 (w\'en) & wen & 一枚铜钱 / cash coin \\
距离 & 里 (l\v{i}) & li & $\approx 556$ 米 / metres \\
距离 & 步 (b\`u) & bu & $\approx 1.54$ 米 / metres \\
\bottomrule
\end{longtable}
\end{center}

\begin{paracol}{2}
\subsection*{数学方法说明}
\addcontentsline{toc}{subsection}{数学方法说明}

本书所涉主要数学方法有四：

\textbf{一、台体体积公式}。长方台之上底为 $a \times b$，下底为 $c \times d$，高为 $h$，则体积
\[V = \frac{h}{3}(ab + cd + ad + bc)\]
此即今之拟柱体公式，秦九韶用以计算土台、堤坝、河渠、仓窖之积。

\textbf{二、衰分均输}。按给定比率分配总量之法。各以本率为列衰，并为法，以总量乘列衰，实如法而一。

\textbf{三、重差之术}。以两次观测（表高与退行距离）推求远处物体之距离与高度，实开三角测量之先声。

\textbf{四、钱陌换算}。以不同陌法之货币统一折算为足陌标准，实为加权平均之应用。
\switchcolumn
\subsection*{Mathematical Methods}
\addcontentsline{toc}{subsection}{Mathematical Methods}

Four principal techniques are employed:

\textbf{1. Frustum volume formula.} For a rectangular frustum with upper $a \times b$, lower $c \times d$, height $h$:
\[V = \frac{h}{3}(ab + cd + ad + bc)\]
Equivalent to the modern prismatoid formula. Used for terraces, dikes, canals, and granaries.

\textbf{2. Proportional distribution (衰分).} Allocation by given ratios. Each party's share is proportional to its stated value.

\textbf{3. Double-difference method (重差術).} Two observations (gnomon height and retreat distance) determine distant heights and distances --- a precursor to trigonometric surveying.

\textbf{4. Currency conversion (陌).} Converting monies of different bundle standards to a uniform full-bundle (足陌) standard --- an application of weighted averages.
\end{paracol}

\vspace{1em}
\begin{paracol}{2}
\subsection*{历史背景}
\addcontentsline{toc}{subsection}{历史背景}

秦九韶于淳佑七年（1247年）撰成《数学九章》，时当南宋末年，国势日蹙。蒙古铁骑压境，故卷八军旅之设尤见苦心。而卷九市物所反映者，乃南宋偏安江左以来商业繁盛、纸币流通、盐茶专卖之经济实况。秦氏以数学济世，其书至今读来犹觉切用。
\switchcolumn
\subsection*{Historical Context}
\addcontentsline{toc}{subsection}{Historical Context}

Qin Jiushao completed the \textit{Shuxue Jiuzhang} in 1247, during the final decades of the Southern Song Dynasty. The military problems in Fascicle 8 directly reflect the desperate strategic situation as the Song faced Mongol invasion. The trade problems in Fascicle 9 reflect the sophisticated commercial economy of the Southern Song, with its paper currency, government salt and tea monopolies, and extensive trade networks.
\end{paracol}

\vspace{2em}
\begin{center}
\longline
\vspace{1em}
{\small\color{sectioncolor}
\textit{--- 卷七至卷九 文言白话对照全译本 完 ---}\\
\textit{--- End of Fascicles 7--9 Bilingual Edition ---}
}
\end{center}

\end{document}
%% ========================================================================
%% END OF DOCUMENT
%% ========================================================================
